%=======================================================================
%  CHAPTER 7 — ADVANCED APPLICATIONS  (3 hrs)
%=======================================================================
\section{Advanced Applications}

{\color{sub}\footnotesize 3 hrs \gap$\cdot$\gap in 22/23 papers \gap$\cdot$\gap 3 topics
\gap$\cdot$\gap 1 numerical (PageRank) $\rightarrow$ Ch 7 companion.}

\lead{Weight of the chapter:}
\begin{itemize}
  \item Web mining \tS{11} \gap $\rightarrow$ \mk{3 hrs of syllabus, asked in 11 papers.
        Best marks-per-hour after Ch 6.}
  \item Time-series data mining \tF{6} \gap $\cdot$\gap Mining object \& multimedia \tP{3}
  \item Mostly \textbf{short notes}, so breadth beats depth here.
\end{itemize}

\hr

\T{7.1 Web Mining}

\Q{What is Web Mining? Briefly explain the structure / taxonomy of Web Mining.}{\m{3+5}\m{6}\m{8}}
{\color{sub}\footnotesize \tS{14} \yr{81 Ba, 80 Bh, 79 Bh, 79 Ba, 78 Bh, 76 Ch, 76 Ash, 75 Ch, 74 Ch, 74 Ash, 73 Ch, 73 Shr, 70 Ch, 69 Ch}}

\sbsr{0.52}{%
\lead{Definition}
\begin{itemize}
  \item \textbf{Web mining} $=$ applying data mining to discover patterns from \textbf{web data}:
        the content, the link structure and the usage logs.
\end{itemize}
\lead{Why the Web is hard \mk{(good opening lines)}}
\begin{itemize}
  \item Huge, growing, and \textbf{widely distributed}.
  \item \textbf{Semi-structured / unstructured} (HTML, not tables).
  \item \textbf{Highly dynamic}; heterogeneous; multilingual.
  \item Only a \textbf{tiny fraction} is relevant to any one user.
  \item Broad topics return hundreds of thousands of documents $\rightarrow$
        \mk{search returns too much, and misses too much}.
\end{itemize}}{%
\lead{The taxonomy \mk{(3 branches --- this IS the answer)}}
\begin{enumerate}
  \item \textbf{Web Content Mining}: mine the \emph{content} of pages (text, images, video).
  \begin{itemize}
    \item \emph{Web page content mining}: extract facts from within pages.
    \item \emph{Search result mining}: summarise / cluster / categorise results.
  \end{itemize}
  \item \textbf{Web Structure Mining}: mine the \emph{hyperlink graph}.
  \begin{itemize}
    \item Links carry authority: a link is a \textbf{vote}.
    \item Algorithms: \textbf{PageRank}, \textbf{HITS}, CLEVER.
    \item Finds authoritative pages, hubs, and web communities.
  \end{itemize}
  \item \textbf{Web Usage Mining}: mine the \emph{server logs} / click-streams.
  \begin{itemize}
    \item \emph{General access pattern tracking}: how the site is used overall.
    \item \emph{Customised usage tracking}: per-user personalisation.
    \item Uses: site restructuring, recommendation, targeted ads, caching, fraud detection.
  \end{itemize}
\end{enumerate}}

\penalty0\vspace{3pt}
\sbsr{0.52}{%
\lead{PageRank \mk{(asked as a short note in 2 papers + 1 numerical)}}
\begin{itemize}
  \item Ranks a page by \textbf{who links to it}, not by what it says.
  \item \textbf{Two ideas}:
  \begin{enumerate}
    \item A link from page $T$ to page $A$ is a \textbf{vote} for $A$.
    \item Votes are \textbf{weighted by the voter's own rank}, and \textbf{split} across all its
          outgoing links.
  \end{enumerate}
  \item \mk{Recursive}: a page is important if important pages point to it.
\end{itemize}
\[ PR(A) = (1-d) + d\sum_{i}\frac{PR(T_i)}{C(T_i)} \]
\begin{itemize}
  \item $T_i$ $=$ pages linking to $A$; \ $C(T_i)$ $=$ number of \textbf{outbound} links from $T_i$.
  \item $d$ $=$ \textbf{damping factor}, usually \textbf{0.85}.
  \item \textbf{Random surfer model}: with probability $d$ the surfer follows a link, with
        probability $(1-d)$ they jump to a random page.
  \item \mk{Why damping is needed}: without it, a group of pages linking only to each other would
        hoard all the rank (a ``rank sink''), and pages with no outlinks would leak it.
  \item Computed \textbf{iteratively} until the values settle; ranks converge quickly.
\end{itemize}}{%
\lead{HITS (for contrast)}
\begin{itemize}
  \item Two scores per page: \textbf{authority} (pointed to by good hubs) and
        \textbf{hub} (points to good authorities).
  \item \mk{Query-dependent} and computed at query time; PageRank is query-independent and
        precomputed.
\end{itemize}}

\vspace{2pt}
\sbs{%
\creamq{Text mining \yr{81 Ba \m{4}}}{}
\begin{itemize}
  \item \textbf{Key steps:}
  \begin{enumerate}
    \item \textbf{Text preprocessing}: tokenise, remove \textbf{stop words}, \textbf{stemming}
          (reduce to root form), case folding.
    \item \textbf{Feature/attribute generation}: build the \textbf{term vector} (bag of words).
    \item \textbf{Feature selection}: drop uninformative terms; weight by \textbf{TF--IDF}.
    \item \textbf{Mining}: classification, clustering, association, topic modelling.
    \item \textbf{Analysis / interpretation} of the results.
  \end{enumerate}
  \item Similarity between documents $\rightarrow$ \textbf{cosine similarity}
        {\footnotesize$\rightarrow$ Ch 2, §2.4}.
  \item \textbf{Information retrieval} measures: precision and recall
        {\footnotesize$\rightarrow$ Ch 3, §3.7}.
\end{itemize}
\lead{Challenges of web mining \yr{75 Ch \m{5}}}
\begin{itemize}
  \item Scale and growth rate; dynamic content; the deep/hidden web.
  \item Noise: ads, navigation bars, boilerplate.
  \item \textbf{Spam}: link farms deliberately gaming the ranking.
  \item Privacy concerns in usage mining; no schema; duplicates and mirrors.
\end{itemize}}{%
\creamq{Mining object and multimedia data}{\m{6}\m{3}\m{4}}
{\color{sub}\footnotesize \tP{3} \yr{75 Ash, 74 Ch, 75 Ch}}
\begin{itemize}
  \item \textbf{Multimedia data} $=$ image, audio, video, text-plus-media in digital form.
  \item \textbf{Feature extraction first} --- you never mine the raw pixels:
  \begin{itemize}
    \item Image: colour histogram, texture, shape, edges.
    \item Audio: pitch, frequency, MFCC. \quad Video: motion vectors, key frames.
  \end{itemize}
  \item \textbf{Tasks}: similarity search (content-based retrieval),
        classification, clustering, association among media features, video-shot detection.
  \item \textbf{Challenges}: very high dimensionality, huge volume, and the
        \mk{semantic gap} --- low-level features do not map cleanly to human meaning.
\end{itemize}
\lead{Visual data mining \yr{75 Ch short note \m{4}}}
\begin{itemize}
  \item Uses \textbf{visualisation} to help the user discover knowledge, combining human perception
        with machine computation.
  \item \textbf{Categories}: data visualisation, mining-\emph{result} visualisation,
        mining-\emph{process} visualisation, and interactive visual mining.
  \item Techniques: scatter plots, box plots, parallel coordinates, pixel-oriented, dendrograms.
  \item \mk{Why it works}: humans spot patterns and outliers in a picture far faster than in a table.
\end{itemize}}

\hr

\T{7.2 Time-Series Data Mining}

\Q{Explain time series data mining with an example.}{\m{6}\m{5}\m{4}\m{3}}
{\color{sub}\footnotesize \tS{11} \yr{82 Bh, 81 Bh, 80 Ba, 79 Ba, 75 Ch, 74 Ch, 73 Ch, 72 Ch, 72 Ka, 71 Shr, 70 Asa} \gap Mostly short notes, but \mk{now asked in 11 papers}.}

\sbs{%
\lead{Definition}
\begin{itemize}
  \item A \textbf{time series} is a sequence of values recorded at \textbf{successive, usually equal,
        time intervals}.
  \item \textbf{Time-series data mining} $=$ finding patterns, trends and regularities in such data.
  \item Eg stock prices, daily temperature, sales per month, sensor readings, ECG, network traffic,
        power consumption.
  \item \mk{What makes it different}: the \textbf{order matters} and consecutive values are
        \textbf{correlated}, so the usual ``independent tuples'' assumption breaks.
\end{itemize}
\lead{The 4 components of a time series \mk{(memorise: T-C-S-I)}}
\begin{itemize}
  \item \textbf{Trend (T)}: the long-term overall direction. Eg steady sales growth.
  \item \textbf{Cyclic (C)}: long, irregular swings, over years. Eg business cycles.
  \item \textbf{Seasonal (S)}: regular, calendar-fixed variation. Eg umbrella sales in monsoon.
  \item \textbf{Irregular / random (I)}: residual noise left after the other three.
  \item Models: multiplicative $Y = T \times C \times S \times I$, or additive $Y = T+C+S+I$.
\end{itemize}}{%
\lead{Tasks}
\begin{itemize}
  \item \textbf{Trend analysis} $+$ \textbf{forecasting} (moving average, exponential smoothing, ARIMA).
  \item \textbf{Similarity search}: find series that behave alike.
  \begin{itemize}
    \item \emph{Whole-sequence} matching vs \emph{subsequence} matching.
    \item Needs \textbf{normalisation} (offset and amplitude) first, else scale dominates.
    \item \textbf{Dynamic Time Warping (DTW)} aligns series that are out of phase or run at
          different speeds --- Euclidean distance cannot. \added
  \end{itemize}
  \item \textbf{Motif discovery}: repeated subsequences. \quad
        \textbf{Anomaly detection}: unusual subsequences {\footnotesize$\rightarrow$ Ch 6}.
  \item \textbf{Segmentation} and \textbf{classification} of series.
\end{itemize}
\lead{Smoothing: moving average}
\begin{itemize}
  \item Replace each value by the mean of a window around it $\rightarrow$ removes the irregular
        component and exposes the trend.
  \item A larger window smooths more but lags more.
\end{itemize}
\lead{Example to quote}
\begin{itemize}
  \item Daily closing prices of a share over 5 years: the \textbf{trend} is the multi-year rise, the
        \textbf{seasonal} part is the recurring festival-quarter bump, the \textbf{cyclic} part is
        the multi-year market swing, and a one-day crash is \textbf{irregular}.
        Mining finds the trend for forecasting and the crash as an anomaly.
\end{itemize}}

\hr

%-----------------------------------------------------------------------
\T{7.3 PYQ Mapping}

{\color{sub}\footnotesize \mk{N} $=$ solved in the Ch 7 Numerical companion.}

\creamq{7.1 Web mining}{}
\begin{enumerate}
  \item \tS{4} What is web mining $+$ its structure / categories / taxonomy.
        \hfill \m{3+5} \yr{80 Bh} \m{6} \yr{74 Ash} \m{8} \yr{76 Ash} \m{5} \yr{75 Ch}
  \item \tP{1} Key steps in \textbf{text mining} $+$ how to find page rank. \hfill \m{4+4} \yr{81 Ba}
  \item \mk{N} \textbf{PageRank} with damping factor 0.85, 5 iterations, initial rank 0.25.
        \hfill \m{8} \yr{76 Ch}
  \item \tP{2} Short notes: Page Rank algorithm \hfill \m{5} \yr{79 Bh} \m{4} \yr{78 Bh}
  \item \tF{3} Short note: Web mining \hfill \m{4} \yr{73 Shr} \m{3} \yr{74 Ch} \yr{70 Ch}
\end{enumerate}

\creamq{7.2 Mining object and multimedia}{}
\begin{enumerate}[start=6]
  \item \tP{1} Explain Web mining and Multimedia mining. \hfill \m{6} \yr{75 Ash}
  \item \tP{2} Short notes: Multimedia mining \hfill \m{3} \yr{74 Ch};
        Visual Data Mining \hfill \m{4} \yr{75 Ch}
\end{enumerate}

\creamq{7.3 Time-series}{}
\begin{enumerate}[start=8]
  \item \tP{1} Explain time series data mining in brief. \hfill \m{6} \yr{72 Ka}
  \item \tP{1} Challenges of web mining $+$ time series data mining w/ example. \hfill \m{5} \yr{75 Ch}
  \item \tF{4} Short note: Time Series Data Mining.
        \hfill \m{5} \yr{80 Ba} \m{3} \yr{74 Ch} \m{4} \yr{72 Ch} \yr{71 Shr}
\end{enumerate}

\creamq{Papers added by the 2026-09 merge}{}
{\color{sub}\footnotesize The six papers of \code{DM4.1BCT.pdf} (2082 Bh, 2081 Bh, 2079 Ba, 2073 Ch, 2070 Asa, 2069 Ch), indexed here for the
first time --- the notes were built before that scan was merged. \mk{N} $=$ solved in this chapter's companion. \added}
\begin{itemize}
  \item Explain \textbf{seasonality} in time-series data. \hfill \m{5} \yr{70 Asa}
  \item Short notes: \textbf{Time Series Data Mining}. \hfill \m{5}\m{4}\m{4}\m{4} \yr{82 Bh, 81 Bh, 79 Ba, 73 Ch}
  \item Short notes: \textbf{Web mining} \yr{79 Ba} $\cdot$ \textbf{www mining} \yr{73 Ch} $\cdot$ \textbf{Page Rank} \yr{69 Ch}. \hfill \m{4}\m{4}\m{5}
  \item Short note: \textbf{Data visualization}. \hfill \m{4} \yr{79 Ba}
\end{itemize}

\lead{Coverage check}
\begin{itemize}
  \item All Ch 7 PYQs covered. 1 numerical \yr{76 Ch} $\rightarrow$ companion.
  \item \mk{This chapter is 3 hrs of syllabus but appears in 16/17 papers}, almost entirely as
        \textbf{4--6 mark short notes}. Learn the \textbf{web mining taxonomy}, the
        \textbf{PageRank formula} and the \textbf{4 time-series components} and you can answer
        nearly every one of them.
  \item \textbf{Gaps filled by research} \added: the \textbf{PageRank formula and the random-surfer
        justification for damping} (your slides name PageRank but never give the equation);
        \textbf{DTW}; the \textbf{T-C-S-I} decomposition; the \textbf{semantic gap} in multimedia mining.
\end{itemize}

\lead{Sources used for this chapter}
\begin{itemize}
  \item \textbf{Han \& Kamber 3rd ed} ch 13 (trends and applications) and the mining-complex-data
        material. \mk{Primary authority.}
  \item \code{Chapter9\_0.ppt} (web mining taxonomy), \code{Chapter9\_1.ppt} (multimedia),
        \code{Chapter9\_2.ppt} (text and web), \code{a1VisualMine.ppt} (visual data mining),
        \code{data mining\_/Chapter 6 \_ 7.pdf}, \textbf{DBK Lectures 18--19}.
\end{itemize}
